\section{Integrasi Lanjut MySQL dengan PHP/API}

Integrasi lanjut MySQL dengan PHP meliputi beberapa pola penting untuk aplikasi production. Pertama, \textbf{connection pooling} atau \textit{persistent connection} (\code{new mysqli('p:localhost', ...)}) menghindari overhead membuka koneksi baru setiap request. Kedua, \textbf{pagination} dari sisi server menggunakan LIMIT dan OFFSET agar tidak memuat seluruh dataset. Ketiga, \textbf{parameterized search} yang menggabungkan beberapa kondisi filter dari form pencarian ke query prepared statement \cite{phpManual}.

\begin{webcode}[language=PHP, caption={Pagination dan pencarian dinamis dengan prepared statement}]
<?php
$page    = max(1, intval($_GET['page'] ?? 1));
$perPage = 10;
$offset  = ($page - 1) * $perPage;
$search  = trim($_GET['q'] ?? '');

// Query dinamis dengan pencarian opsional
$sql    = "SELECT * FROM produk";
$params = [];
$types  = "";

if (!empty($search)) {
  $sql    .= " WHERE nama LIKE ?";
  $params[] = "%{$search}%";
  $types   .= "s";
}
$sql .= " ORDER BY id DESC LIMIT ?, ?";
$params[] = $offset;
$params[] = $perPage;
$types   .= "ii";

$stmt = $conn->prepare($sql);
$stmt->bind_param($types, ...$params);
$stmt->execute();
$result = $stmt->get_result();
$produk = $result->fetch_all(MYSQLI_ASSOC);
$stmt->close();

// Hitung total halaman
$countSql = "SELECT COUNT(*) AS total FROM produk";
if (!empty($search)) {
  $countSql .= " WHERE nama LIKE ?";
}
// ... fetch total untuk pagination links
?>
\end{webcode}

